\centerline{\bf \Huge ММО 2025}
\centerline{\bf \Huge 9 класс}

\begin{flushright}
        \scriptsize{}
\end{flushright}

\problem{1}
Положительные числа $a$ и $b$ таковы, что $a-b=a / b$. Что больше, $a+b$ или $a b$ ?

\problem{2}
В треугольнике $A B C$ на продолжении медианы $C M$ за точку $C$ отметили точку $K$ так, что $A M=C K$. Известно, что угол $B M C$ равен $60^{\circ}$. Докажите, что $A C=B K$.

\problem{3}
В узлах сетки клетчатого прямоугольника $4 \times 5$ расположены 30 лампочек, изначально все они погашены. За ход разрешается провести любую прямую, не задевающую лампочек (размерами лампочек следует пренебречь, считая их точками), такую, что с какой-то одной стороны от неё ни одна лампочка не горит, и зажечь все лампочки по эту сторону от прямой. Каждым ходом нужно зажигать хотя бы одну лампочку. Можно ли зажечь все лампочки ровно за четыре хода?

\problem{4}
Точки $O$ и $I$ - центры описанной и вписанной окружностей неравнобедренного треугольника $A B C$. Две равные окружности касаются сторон $A B, B C$ и $A C, B C$ соответственно; кроме этого, они касаются друг друга в точке $K$. Оказалось, что $K$ лежит на прямой $O I$. Найдите $\angle B A C$.

\problem{5}
Рациональные числа $x, y$ и $z$ таковы, что все числа $x+y^2+z^2$, $x^2+y+z^2$ и $x^2+y^2+z$ целые. Верно ли, что число $2 x$ целое?

\problem{6}
Пусть $p$ и $q$ - взаимно простые натуральные числа. Лягушка прыгает по числовой прямой, начиная в точке $0$, каждый раз либо на $p$ вправо, либо на $q$ влево. Однажды лягушка вернулась в $0$. Докажите, что для любого натурального $d<p+q$ найдутся два числа, посещённые лягушкой и отличающиеся на $d$.